\documentclass[aps,prx,twocolumn,superscriptaddress,nofootinbib]{revtex4-2}

\usepackage{amsmath,amssymb}
\usepackage{booktabs}
\usepackage{hyperref}
\hypersetup{colorlinks=true,linkcolor=blue,citecolor=blue,urlcolor=blue}

\begin{document}

\title{Held-out interpolation does not validate zero-noise extrapolation:
a $513\times$ amplification of pure shot noise on a noiseless model}

\author{Venkata Rao Allu}
\email{alluvenkat11@gmail.com}
\affiliation{Independent researcher}

\date{\today}

\begin{abstract}
Zero-noise extrapolation fits a noise-scaling curve to measurements taken
at amplified noise and evaluates it at zero. Implementations are commonly
validated by holding out one amplified-noise point and checking the fit
predicts it. We show this validation is systematically blind to the failure
mode that matters, and quantify the gap. In a per-Pauli extrapolation over
$756$ curves, the held-out prediction error was $7.5\times10^{-4}$ on
quantities bounded to $[-1,1]$ --- an excellent fit by any conventional
standard --- while the same fits, evaluated at the extrapolation target,
turned a $0.0652$~kcal/mol energy error into $33.48$~kcal/mol, an
amplification of $513\times$. The input carried \emph{no} device noise: it
was multinomial shot noise drawn from exact statevector probabilities for a
noiseless model, so the entire discrepancy was manufactured by the
estimator. We trace the mechanism to the variance inflation factor of the
fit evaluated off its data support. At the held-out point this factor is
below unity for three of the four model classes used; at the extrapolation
target it reaches $4.0$, $6.9\times10^{7}$ and $1.7\times10^{20}$ for the
quadratic, exponential and rational classes. Shot noise then drives
model-class selection toward the flexible nonlinear families --- chosen for
$62\%$ of curves --- precisely because they interpolate best, and those are
the classes whose extrapolation is worst conditioned. A layer-by-layer
control shows circuits, measurement, reconstruction and the extrapolator's
own algebra are each exact to $10^{-11}$ or better on noiseless input,
isolating the pathology to the interaction of finite sampling with
off-support evaluation. We recommend that any zero-noise extrapolation
report its variance inflation at the evaluation point, and that validation
be performed by extrapolating to a withheld \emph{lower}-noise point rather
than by interpolation.
\end{abstract}

\maketitle

\section{Introduction}

Zero-noise extrapolation (ZNE) is the most widely deployed error
mitigation technique for near-term quantum
processors~\cite{temme2017,li2017,kandala2019}. Its logic is simple:
deliberately amplify the noise in a circuit by a controlled factor
$\lambda$, measure an observable at several $\lambda>1$, fit a curve, and
evaluate it at $\lambda=0$.

The practical question is how one knows the fit is trustworthy. The
natural answer, and the one we adopted, is cross-validation: withhold one
of the measured amplification levels, fit to the rest, and check that the
fit predicts the withheld point. If it does, the functional form is
presumably adequate.

This paper reports that the check is close to worthless for its intended
purpose, and explains why in a way we believe generalises. The withheld
point lies \emph{between} or at the edge of the remaining data; the
evaluation point $\lambda=0$ lies outside it, on the opposite side from
every measurement. These are different statistical problems. A fit can be
excellent at the former and catastrophically ill-conditioned at the
latter, and the conventional validation cannot see the difference by
construction.

We did not discover this by analysis. We discovered it because a
production pipeline produced an energy error of $33.48$~kcal/mol --- some
$500$ times the uncorrected error --- on an input that contained no
hardware noise whatsoever, and we were forced to find out why.

\section{The pipeline and the control experiment}

The observable is the ground-state energy of H$_4$ (linear chain,
$1.0$~\AA{}, STO-3G, $4$ electrons in $4$ orbitals), computed through an
entanglement-forged decomposition~\cite{eddins2022} into two $4$-qubit
registers. Noise is amplified by native-gate folding, $G \to
G\,(G^{-1}G)^{k}$, applied to the two-qubit entangling gate, giving
amplification levels $\lambda \in \{1,3,5,7,9\}$. The energy requires
$756$ Pauli expectation values across all preparation slots, and the
implementation fits an independent noise-scaling curve to each.

Model selection is itself cross-validated, in two stages: a model class is
chosen per curve from $\{$linear, quadratic, exponential, rational$\}$ by
held-out prediction error, then refit. This is more careful than a fixed
functional form, and it is the specific reason the failure is as large as
it is.

Before attributing the blowup to the extrapolator we eliminated every
other component. Table~\ref{tab:control} reports a layer-by-layer control
on exact statevector input, where the correct answer is known in closed
form at every amplification level.

\begin{table}[t]
\caption{Layer-by-layer control on noiseless statevector input. Each layer
adds one stage of the production pipeline. Values are deviations from the
exact result; folding is an identity operation on noiseless input, so the
correct behaviour is independence from $\lambda$.}
\label{tab:control}
\begin{ruledtabular}
\begin{tabular}{llc}
Layer & What it tests & Deviation \\
\colrule
A & circuit construction, folding identity & $\le1.1\times10^{-14}$ \\
B & Pauli expectation extraction & $\le1.3\times10^{-15}$ \\
C1 & per-slot state reconstruction & $\le2.0\times10^{-14}$ \\
C2 & energy assembly (kcal/mol) & $\le5.5\times10^{-11}$ \\
D & the extrapolator's own algebra & $6.1\times10^{-12}$ \\
\end{tabular}
\end{ruledtabular}
\end{table}

Every layer is exact. Folding at $\lambda=9$ reproduces $\lambda=1$ to
$10^{-14}$; the reconstructed energy is correct to $5\times10^{-11}$~kcal/mol
at every level; and the two-stage extrapolator, fed exact data, selects
the linear class with a held-out error of $4\times10^{-16}$ and returns
the exact answer. There is no bug in any of it.

\section{The failure}

We then changed exactly one thing. Instead of exact probabilities we drew
$100{,}000$ multinomial samples per circuit from those same exact
probabilities, over $8$ seeds --- statistically identical to a real
submission to a noiseless simulator, and containing no device noise of any
kind. The identical production code was run verbatim.

\begin{table}[t]
\caption{Pure shot noise through the production extrapolator, $756$
curves, $100{,}000$ shots, noiseless model. The input contains no device
noise, so the correct $\lambda=0$ answer equals the $\lambda=1$ value.}
\label{tab:blowup}
\begin{ruledtabular}
\begin{tabular}{lc}
Quantity & Value \\
\colrule
Stage-1 mean held-out error & $8.19\times10^{-4}$ \\
Stage-2 mean held-out error & $7.50\times10^{-4}$ \\
\colrule
Energy error at $\lambda=1$ (raw) & $0.0652$~kcal/mol \\
Energy error at $\lambda=0$ (extrapolated) & $\mathbf{33.48}$~kcal/mol \\
Amplification & $\mathbf{513\times}$ \\
\colrule
Curves rejected for $|\langle P\rangle|>1$ & $115/756$ $(15.2\%)$ \\
Worst rejected extrapolations & $17{,}295$ and $-8{,}214$ \\
\colrule
Class chosen: rational & $207/756$ \\
Class chosen: exponential & $262/756$ \\
Class chosen: linear & $143/756$ \\
Class chosen: quadratic & $144/756$ \\
\end{tabular}
\end{ruledtabular}
\end{table}

Table~\ref{tab:blowup} is the whole paper in one place. The held-out error
is $7.5\times10^{-4}$ on quantities bounded to $[-1,1]$: by the standard
diagnostic, these are good fits. The same fits evaluated at $\lambda=0$
produce an energy error $513$ times the unmitigated one, and every bit of
that $33.4$~kcal/mol was created by the estimator, because the input had
no device noise to remove.

Two secondary numbers show how violently the fits misbehave off support.
$115$ of $756$ curves extrapolated outside $[-1,1]$ and had to be
rejected --- a Pauli expectation value is bounded by $\pm1$ by
construction --- and the worst two returned $17{,}295$ and $-8{,}214$.
A fit that predicts a held-out point to four decimal places and then
returns five orders of magnitude outside a hard physical bound is not
mildly misspecified; it is being evaluated somewhere it carries no
information.

\section{Mechanism: variance inflation off the data support}

The diagnosis is the conditioning of the fit at its evaluation point. For
a linear-in-parameters fit with design matrix $X$ over the measured
amplification levels, the variance of the prediction at a point $x$ is
\begin{equation}
\mathrm{Var}[\hat y(x)] = \sigma^2\, x^\top (X^\top X)^{-1} x
\equiv \sigma^2\, \mathrm{VIF}(x),
\end{equation}
where $\sigma^2$ is the per-point measurement variance and
$\mathrm{VIF}(x)$ is a variance inflation factor determined entirely by
the design --- the choice of amplification levels and functional form ---
and not at all by the data. It can be computed before any measurement is
taken.

Table~\ref{tab:vif} evaluates it at the held-out validation point
($\lambda=9$) and at the extrapolation target ($\lambda=0$), for the four
classes in use, with the nonlinear classes linearised about their fitted
parameters.

\begin{table}[t]
\caption{Variance inflation at the validation point versus the evaluation
point, for the design $\lambda\in\{1,3,5,7,9\}$. For the nonlinear classes
the statistics are over the $28$ representative curves fitted; means and
medians are both given because the distributions are heavy-tailed.}
\label{tab:vif}
\begin{ruledtabular}
\begin{tabular}{lccc}
Class & $\mathrm{VIF}(\lambda{=}9)$ & $\mathrm{VIF}(\lambda{=}0)$ & ratio \\
\colrule
linear      & $0.600$ & $0.825$ & $1.4$ \\
quadratic   & $0.886$ & $2.115$ & $2.4$ \\
exponential (median) & $0.577$ & $4.008$ & $6.9$ \\
exponential (mean)   & $0.516$ & $6.94\times10^{7}$ & $1.3\times10^{8}$ \\
rational (median)    & $0.460$ & $903.5$ & $2.0\times10^{3}$ \\
rational (mean)      & $2385$  & $1.71\times10^{20}$ & $7.2\times10^{16}$ \\
\end{tabular}
\end{ruledtabular}
\end{table}

The structure is unambiguous. At the held-out point the inflation factor
is below unity for every class but one: the measurement is being
\emph{averaged}, and the validation looks good because interpolation
genuinely is a well-posed problem on this design. At the evaluation point
the same designs inflate the variance by factors reaching $10^{20}$.

The two columns are not measuring the same thing, and no amount of
agreement in the first constrains the second. This is the precise sense in
which held-out interpolation fails to validate extrapolation: the
diagnostic is evaluated where the problem is well-conditioned and the
answer is taken where it is not.

\subsection{Why shot noise makes it worse, not merely noisier}

A subtler effect compounds this, and it explains the factor of $513$
rather than a factor of a few. Model selection is driven by held-out
\emph{interpolation} error. The flexible three-parameter nonlinear classes
interpolate best --- they have more freedom to accommodate the particular
realisation of shot noise in the fitted points --- so selection prefers
them: rational and exponential were chosen for $469$ of $756$ curves,
$62\%$ of the total. Those are exactly the two classes whose
$\mathrm{VIF}(\lambda{=}0)$ is astronomical.

Finite sampling therefore does not simply add variance to the
extrapolation. It actively steers model selection toward the classes whose
extrapolation is worst conditioned, because the criterion used to select
them rewards the behaviour that is orthogonal to the one that matters. The
more shot noise, the stronger the preference. This is a selection
pathology, not an estimation one, and increasing the shot count mitigates
it only slowly.

\subsection{The pathology is in the fitting functions, not the curve count}

One natural objection is that fitting $756$ independent curves is simply
reckless, and that a parsimonious estimator would be safe. We tested this.
H$_4$'s forged state at Schmidt rank $K=6$ lives in a six-dimensional
Hamming-weight-two sector with real amplitudes and unit norm, leaving five
free parameters, against $35$ for a general density matrix on the same
register. We built that five-parameter physical-manifold estimator and it
is a substantial improvement at a \emph{single} amplification level:
$3.4\times$ better on one device noise model and $11\times$ better on the
noiseless one, and it passes an adversarial control --- fitted to shuffled
or random data it returns errors of order $10^{3}$~kcal/mol, so it is not
covertly injecting the answer.

Extrapolating its five components across $\lambda$ reproduced the same
blowup. Reducing $756$ curves to $5$ did not help, because the problem was
never the number of curves; it is that off-support evaluation of a fitted
nonlinear function is ill-conditioned regardless of how few such functions
there are.

\section{What does transfer}

The constructive counterpart is worth stating, because the conclusion is
not that error mitigation is hopeless. Across this project every method
that applies a \emph{direct} correction --- a forward model of the noise
channel inverted at the measured amplification level, in the manner of
probabilistic error
cancellation~\cite{temme2017,vandenberg2023} --- behaved stably, and every
method that fits a curve and evaluates it off support failed. The
distinguishing feature is not sophistication; it is whether the estimator
is evaluated inside or outside the support of its data.

Direct corrections fail differently, and better. In a separate
pre-registered experiment on a trapped-ion device, an analytic direct
correction of this kind missed its target by a factor of $4.2$ and was
found, by refitting a single scalar blind to the reference energy, to
under-correct by a factor of $\approx1.8$~\cite{companion}. That is a
bounded, diagnosable, one-parameter discrepancy. It is a different species
of failure from manufacturing $33$~kcal/mol out of shot noise.

\section{Recommendations}

\begin{enumerate}
\item \textbf{Report the variance inflation at the evaluation point.}
$\mathrm{VIF}(\lambda{=}0)$ depends only on the amplification levels and
the functional form, so it can and should be computed and published
before any data is taken. A design with
$\mathrm{VIF}(0)\gg\mathrm{VIF}(\lambda_{\mathrm{val}})$ is announcing in
advance that its validation is uninformative.

\item \textbf{Validate by extrapolation, not interpolation.} Withhold the
\emph{lowest} amplification level and extrapolate to it from the higher
ones. This tests the operation actually performed. It is strictly
available whenever a conventional held-out check is, and it costs nothing
extra.

\item \textbf{Enforce physical bounds as a hard gate, not a filter.} That
$15\%$ of curves extrapolated outside $[-1,1]$, two of them by four
orders of magnitude, was recorded as a rejection count rather than treated
as a failure of the procedure. A procedure that violates a construction
bound on one curve in seven has been falsified, and discarding the
violations conceals that.

\item \textbf{Do not select model class by interpolation error} when the
quantity of interest is an extrapolation. The criterion rewards the
classes that extrapolate worst.

\item \textbf{Prefer direct correction where a forward noise model is
available.} The amplification-and-extrapolation route trades a bounded
modelling error for an unbounded conditioning one.
\end{enumerate}

\section{Limitations}

These results come from one observable, one molecule, one amplification
design ($\lambda\in\{1,3,5,7,9\}$ by two-qubit gate folding) and one
implementation. The VIF argument is general --- it is a property of linear
least squares --- but the specific magnitudes in Table~\ref{tab:vif}
belong to this design, and a different set of amplification levels would
give different numbers. We make no claim that published ZNE results are
wrong; many use fixed low-order functional forms, for which
Table~\ref{tab:vif} shows the inflation to be modest ($1.4\times$ linear,
$2.4\times$ quadratic). Our claim is narrower and, we think, unavoidable:
the held-out interpolation check does not bound the extrapolation error,
the gap can be twenty orders of magnitude, and shot-noise-driven model
selection moves estimators toward the regime where it is largest.

\section{Conclusion}

A per-Pauli zero-noise extrapolation with excellent held-out
interpolation error, $7.5\times10^{-4}$, amplified pure shot noise into a
$33.48$~kcal/mol energy error on a model containing no device noise at
all, a factor of $513$ above the unmitigated value. Every other component
of the pipeline is exact to $10^{-11}$ or better. The mechanism is the
variance inflation of a fitted curve evaluated outside the support of its
data, which reaches $10^{20}$ for the flexible nonlinear classes that
shot-noise-driven selection prefers precisely because they interpolate
well.

Held-out interpolation and zero-noise extrapolation are different
statistical operations, and success at the first is not evidence about the
second. We recommend reporting $\mathrm{VIF}$ at the evaluation point and
validating by extrapolation to a withheld lower-noise point, both of which
are free.

\begin{acknowledgments}
This work was supported by the IonQ Research Credits Program. The author
thanks Vadim Karpusenko (IonQ) for technical guidance, and in particular
for the native-gate diagnosis that made the folding control of
Table~\ref{tab:control} possible.

In accordance with arXiv's policy on generative AI tools, we disclose that
a large language model (Anthropic Claude) was used as a coding and
drafting assistant during this work: for software implementation,
analysis scripting, and preparation of this manuscript. All scientific
claims, experimental design decisions and interpretations are the author's
own, and the author takes full responsibility for the content. All
numerical results were produced by the committed code and are
reproducible from the public repository.
\end{acknowledgments}

\section*{Data and code availability}

All code, result files and a complete iteration-by-iteration research log
retaining every negative result are at
\url{https://github.com/Venkatallu11/quantum-chemistry-vqe}. The control
experiment of Table~\ref{tab:control} is
\texttt{vqe/task29a\_ideal\_control\_diagnosis.py}; the shot-noise blowup
of Table~\ref{tab:blowup} is
\texttt{vqe/task29a\_layerD2\_realistic.py}; the conditioning analysis of
Table~\ref{tab:vif} is
\texttt{vqe/task29g\_extrapolation\_conditioning.py}. Each writes the
corresponding \texttt{*\_results.json} from which the tables are taken
directly.

\begin{thebibliography}{9}

\bibitem{temme2017} K. Temme, S. Bravyi, and J. M. Gambetta,
\emph{Error mitigation for short-depth quantum circuits},
Phys. Rev. Lett. \textbf{119}, 180509 (2017).

\bibitem{li2017} Y. Li and S. C. Benjamin, \emph{Efficient variational
quantum simulator incorporating active error minimization},
Phys. Rev. X \textbf{7}, 021050 (2017).

\bibitem{kandala2019} A. Kandala, K. Temme, A. D. C\'orcoles,
A. Mezzacapo, J. M. Chow, and J. M. Gambetta, \emph{Error mitigation
extends the computational reach of a noisy quantum processor},
Nature \textbf{567}, 491 (2019).

\bibitem{eddins2022} A. Eddins, M. Motta, T. P. Gujarati, S. Bravyi,
A. Mezzacapo, C. Hadfield, and S. Sheldon, \emph{Doubling the size of
quantum simulators by entanglement forging}, PRX Quantum \textbf{3},
010309 (2022).

\bibitem{vandenberg2023} E. van den Berg, Z. K. Minev, A. Kandala, and
K. Temme, \emph{Probabilistic error cancellation with sparse
Pauli--Lindblad models on noisy quantum processors},
Nat. Phys. \textbf{19}, 1116 (2023).

\bibitem{companion} V. R. Allu, \emph{A pre-registered
entanglement-forged H$_4$ benchmark on trapped-ion hardware: accurate
noise modelling, incomplete error correction}, companion manuscript
(2026).

\end{thebibliography}

\end{document}
